\chapter{Documentarea}

\section{Stadiul actual al domeniului}
Jocul de Sudoku standard constă într-o grilă de dimensiune $9 \times 9$, împărțită în 9 sub-grile (regiuni sau blocuri) de dimensiune $3 \times 3$. Obiectivul jocului este completarea celulelor libere cu cifre de la 1 la 9, astfel încât fiecare linie, fiecare coloană și fiecare sub-grilă $3 \times 3$ să conțină fiecare cifră de la 1 la 9 o singură dată.

Din punct de vedere matematic și algoritmic, Sudoku este o problemă clasică de satisfacere a constrângerilor (\textit{Constraint Satisfaction Problem} -- CSP) \cite{russell2020}. În forma sa generalizată pe o grilă de dimensiune $n^2 \times n^2$, problema a fost demonstrată ca fiind \textbf{NP-completă} \cite{yato2003}. Pentru grila clasică $9 \times 9$, numărul total de configurații complete valide a fost determinat în anul 2005 de către Bertram Felgenhauer și Frazer Jarvis ca fiind egal cu $6.670.903.752.021.072.936.960 \approx 6,67 \times 10^{21}$ \cite{felgenhauer2006}. De asemenea, în anul 2012, Gary McGuire și colaboratorii săi au demonstrat că numărul minim de indicii (\textit{clues}) necesare pentru ca o grilă de Sudoku să admită o soluție unică este 17 \cite{mcguire2014}.

În stadiul actual, metodele de rezolvare și generare automată se împart în următoarele direcții principale:
\begin{itemize}
    \item \textbf{Căutarea cu revenire (Backtracking / DFS):} Abordarea fundamentală care parcurge celulele libere și testează secvențial valorile valide din domeniul $\{1, \dots, 9\}$. Dacă pentru o celulă curentă nu există nicio valoare validă, algoritmul revine la celula anterioară (\textit{backtrack}) și încearcă următoarea opțiune.
    \item \textbf{Propagarea constrângerilor și euristici CSP:} Optimizează căutarea prin reducerea domeniilor de valori posibile înainte de ramificare \cite{norvig2006}. Se utilizează euristici precum \textit{Minimum Remaining Values} (MRV -- selectarea celulei cu cei mai puțini candidați posibili) și reguli de eliminare logică (\textit{Naked Singles}, \textit{Hidden Singles}).
    \item \textbf{Algoritmul X și Dancing Links (DLX):} Formulat de Donald E. Knuth \cite{knuth2000}, transformă rezolvarea unei grile de Sudoku într-o problemă de acoperire exactă (\textit{Exact Cover Problem}) pe o matrice binară cu 729 de linii (câte 9 candidați pentru fiecare dintre cele 81 de celule) și 324 de coloane (cele 4 tipuri de constrângeri: poziție, linie, coloană, regiune). Este una dintre cele mai rapide metode pentru verificarea unicității soluției.
    \item \textbf{Algoritmi de generare a grilelor:} Generarea unui puzzle valid cu soluție unică se realizează în două faze:
    \begin{enumerate}
        \item \textit{Construirea unei grile complete valide:} Se populează aleator cele trei sub-grile diagonale $3 \times 3$ (deoarece sunt independente între ele), apoi se aplică algoritmul de rezolvare cu selecție randomizată a candidaților pentru a umple restul grilei.
        \item \textit{Eliminarea controlată a valorilor (Digging Holes):} Se elimină pe rând cifre din grila completă în ordine aleatoare (sau simetrică). După fiecare eliminare, se apelează modulul de rezolvare pentru a număra soluțiile posibile; dacă grila admite mai mult de o soluție, cifra eliminată este restaurată. Numărul de celule rămase și tehnicile necesare rezolvării determină nivelul de dificultate (Ușor, Mediu, Dificil).
    \end{enumerate}
\end{itemize}

\section{Surse de documentare}
Documentarea teoretică și tehnică se bazează pe următoarele categorii de surse (detaliate în secțiunea Bibliografie):
\begin{itemize}
    \item \textbf{Articole științifice și combinatorică:} Lucrarea lui Donald Knuth privind structurile \textit{Dancing Links} \cite{knuth2000}, analiza matematică a grilelor de Sudoku realizată de Felgenhauer și Jarvis \cite{felgenhauer2006}, demonstrația de NP-completitudine a lui Yato și Seta \cite{yato2003} și studiul privind numărul minim de indicii realizat de McGuire et al. \cite{mcguire2014}.
    \item \textbf{Literatură de specialitate în Inteligență Artificială:} Tratatul lui Stuart Russell și Peter Norvig \cite{russell2020}, pentru formalizarea problemelor de satisfacere a constrângerilor (CSP), căutarea Backtracking și euristicile de tip MRV.
    \item \textbf{Studii de caz și implementări practice:} Articolul tehnic al lui Peter Norvig privind propagarea constrângerilor în Sudoku \cite{norvig2006}.
\end{itemize}

\section{Planul individual de lucru}
Pentru realizarea proiectului de practică, activitatea a fost împărțită în următoarele etape de lucru:
\begin{enumerate}
    \item \textbf{Etapa 1 (S1-E1) -- Definirea temei:} Configurarea depozitului Git, inițializarea caietului de practică, stabilirea scopului, a obiectivelor și a contextului aplicației.
    \item \textbf{Etapa 2 (S1-E2) -- Documentarea:} Studiul algoritmilor existenți pentru generarea și rezolvarea jocurilor de tip Sudoku (Backtracking, CSP, Exact Cover), selectarea surselor bibliografice și stabilirea planului individual de lucru.
    \item \textbf{Etapa 3 (S1-E3) -- Analiza și Proiectarea:} Definirea cerințelor funcționale și nefuncționale, alegerea limbajului de programare și a tehnologiilor, proiectarea arhitecturii aplicației, a structurilor de date pentru reprezentarea grilei $9 \times 9$ și a diagramelor UML.
    \item \textbf{Etapa 4 (S1-E4) -- Realizarea (Implementarea):} 
    \begin{itemize}
        \item Implementarea validatorului de reguli și a algoritmului de rezolvare (Solver) capabil să găsească soluția și să verifice unicitatea acesteia.
        \item Implementarea generatorului de grile valide (Generator) pe mai multe niveluri de dificultate.
        \item Dezvoltarea interfeței cu utilizatorul pentru jucarea efectivă, verificarea mutărilor și solicitarea de indicii.
    \end{itemize}
    \item \textbf{Etapa 5 (S1-E5) -- Testarea și Concluziile:} Verificarea corectitudinii generării (unicitatea soluției), măsurarea timpilor de execuție pentru diferite grade de dificultate, validarea funcționalităților și redactarea concluziilor finale ale stagiului de practică.
\end{enumerate}